\element{1.3.1}{
  \begin{questionmult}{OscillateurWienComposition}\bareme{haut=2,p=0}
    De quoi est constitué un oscillateur de Wien ?
    \begin{multicols}{2}
      \begin{reponses}
        \bonne{Un amplificateur}
        \bonne{Un filtre passe bande}
        \mauvaise{Un comparateur à hystérésis}
        \mauvaise{Un intégrateur}
      \end{reponses}
    \end{multicols}
  \end{questionmult}
}
\element{1.3.1}{
  \begin{questionmult}{OscillateurRelaxationComposition}\bareme{haut=2,p=0}
    De quoi est constitué un oscillateur à relaxation ?
    \begin{multicols}{2}
      \begin{reponses}
        \mauvaise{Un amplificateur}
        \mauvaise{Un filtre passe bande}
        \bonne{Un comparateur à hystérésis}
        \bonne{Un intégrateur}
      \end{reponses}
    \end{multicols}
  \end{questionmult}
}
\setgroupmode{1.3.1}{cyclic}

\element{1.3.2}{
  \begin{questionmult}{WienDemarrage}\bareme{haut=2,p=0}
    Pour que des oscillations démarrent dans un oscillateur quasi-sinusoïdal, il faut que
    \begin{multicols}{2}
      \begin{reponses}
        \mauvaise{L'ALI soit dans un montage instable}
        \mauvaise{Le produit des fonctions de transfert de chaque système soit égal à $1$}
        \mauvaise{Chaque élément du système bouclé soient instable}
        \bonne{Le système bouclé soit instable}
      \end{reponses}
    \end{multicols}
  \end{questionmult}
}
\element{1.3.2}{
  \begin{questionmult}{WienExistance}\bareme{haut=2,p=0}
    Pour que les oscillations sinusoïdales existent dans un oscillateur quasi-sinusoïdal, il faut que
    \begin{multicols}{2}
      \begin{reponses}
        \mauvaise{L'ALI soit dans un montage instable}
        \bonne{Le produit des fonctions de transfert de chaque système soit égal à $1$}
        \mauvaise{Chaque élément du système bouclé soient instable}
        \mauvaise{Le système bouclé soit instable}
      \end{reponses}
    \end{multicols}
  \end{questionmult}
}
\setgroupmode{1.3.2}{cyclic}

\element{1.3.3}{
  \begin{questionmult}{WienFrequenceDepend}\bareme{haut=2,p=0}
    La fréquence des oscillations de l'oscillateur de Wien dépend de
    \begin{multicols}{2}
      \begin{reponses}
        \bonne{$R$}
        \bonne{$C$}
        \mauvaise{$R_1$ et $R_2$}
        \mauvaise{$V_{\text{sat}}$}
      \end{reponses}
    \end{multicols}
  \end{questionmult}
}
\element{1.3.3}{
  \begin{questionmult}{RelaxationPeriodeDepend}\bareme{haut=2,p=0}
    La période des oscillations de l'oscillateur à relaxation dépend de
    \begin{multicols}{2}
      \begin{reponses}
        \bonne{$R$ et $C$}
        \bonne{$R_1$ et $R_2$}
        \mauvaise{$V_{\text{sat}}$}
        \mauvaise{Aucune de ces grandeurs}
      \end{reponses}
    \end{multicols}
  \end{questionmult}
}
\setgroupmode{1.3.3}{cyclic}

\element{1.3.5}{
  \begin{question}{RelaxationAmplitude}\bareme{1}
    Dans l'oscillateur à relaxation, l'amplitude du signal créneau vaut
    \begin{multicols}{2}
      \begin{reponses}
        \mauvaise{$\frac{R_1}{R_2}V_{\text{sat}}$}
        \mauvaise{$\frac{R_2}{R_1}V_{\text{sat}}$}
        \mauvaise{$\frac{R_1}{R_1+R_2}V_{\text{sat}}$}
        \bonne{$V_{\text{sat}}$}
      \end{reponses}
    \end{multicols}
  \end{question}
}
\element{1.3.5}{
  \begin{question}{WienAmplitude}\bareme{1}
    Dans l'oscillateur de Wien, l'amplitude des oscillations est fixée par
    \begin{multicols}{2}
      \begin{reponses}
        \bonne{La saturation de l'ALI}
        \mauvaise{Le rapport $R_2/R_1$}
        \mauvaise{Les conditions initiales}
        \mauvaise{Le produit $RC$}
      \end{reponses}
    \end{multicols}
  \end{question}
}
\setgroupmode{1.3.5}{cyclic}

\element{1.3.6}{
  \begin{questionmult}{WienHarmoniques}\bareme{haut=2,p=0}
    Dans l'oscillateur de Wien, si on augmente $R_2$ au-delà de $2R_1$,
    \begin{reponses}
      \bonne{Le spectre des tensions comporte plus d'harmoniques}
      \bonne{La sortie du filtre reste plus sinusoïdale que celle de l'amplificateur}
      \mauvaise{L'amplitude des oscillations augmente indéfiniment}
      \mauvaise{Les oscillations s'arrêtent}
    \end{reponses}
  \end{questionmult}
}
\element{1.3.6}{
  \begin{questionmult}{RelaxationR2}\bareme{haut=2,p=0}
    Dans l'oscillateur à relaxation, si on augmente $R_2$ (en gardant $R_2>R_1$),
    \begin{reponses}
      \bonne{La période diminue}
      \bonne{L'amplitude du signal triangulaire diminue}
      \mauvaise{L'amplitude du signal créneau diminue}
      \mauvaise{Les oscillations s'arrêtent}
    \end{reponses}
  \end{questionmult}
}
\setgroupmode{1.3.6}{cyclic}
